\documentclass{article}

\title{Notes on The C Programming Language}
\author{Nicoll Douglas}

\usepackage{float}
\usepackage[T1]{fontenc}
\usepackage{tablefootnote}
\usepackage{listings}
\usepackage{xcolor}
\usepackage[expansion=false]{microtype}
\usepackage{hyperref}

\hypersetup{
    hidelinks
}

\DisableLigatures[>,<]{encoding = T1, family = tt* }

\lstset{
    language=C,
    basicstyle=\ttfamily\small,
    keywordstyle=\color{blue},
    commentstyle=\color{gray},
    numbers=left,
    breaklines=true
}

\newcommand{\code}[1]{\texttt{#1}}

\begin{document}

\maketitle
\tableofcontents

\section{Types, Operators \& Expressions}

\subsection{Variable Names}

\begin{itemize}
    \item Must only contain letters, digits, and underscores.
    \item Must start with a letter or underscore\footnote{Library functions often start with underscores so avoid.}.
\end{itemize}

\subsection{Data Types and Sizes}

\begin{table}[H]
    \centering
    \begin{tabular}{|l|l|l|}
        \hline
        \textbf{Type} & \textbf{Size (bytes)} & \textbf{Description} \\
        \hline
        \code{char} & 1 & Holds one character \\
        \code{short} / \code{short int} & 2\tablefootnote{C standard minimum 2 bytes, typically 2 bytes.} & Smaller integer \\
        \code{int} & 4\tablefootnote{C standard minimum 2 bytes and at least as big as \code{short}; typically 4 bytes.} & Natural integer size on host machine \\
        \code{long} / \code{long int} & 8\tablefootnote{C standard minimum 4 bytes and at least as big as \code{int}; typically 8 bytes.} & Bigger integer \\
        \code{float} & 4 & Single-precision floating point \\
        \code{double} & 8 & Double-precision floating point \\
        \code{long double} & 16\tablefootnote{C standard at least as long as \code{double}; usually 8, 12 or 16 bytes.} & Extended-precision floating-point \\
        \hline
    \end{tabular}
    \caption{Type, size, and description of primitive data types in C}
\end{table}

\begin{itemize}
    \item The \code{signed} or \code{unsigned} qualifiers may be applied to \code{char} or any integer type.
    \item \code{unsigned} integers are positive or 0.
    \item \code{unsigned} integers obey the laws of arithmetic module $2^n$ where $n$ is the number of bits in the data type.
    \item \code{int}, \code{short}, and \code{long} are all signed by default.
    \item \code{char} can be \code{unsigned} or \code{signed} by default (not mandated in C standard) but is usually \code{signed}.
\end{itemize}

\subsection{Constants}

\subsubsection{Integer Constants}

\paragraph{Suffixes} Integer constants may contain the given suffixes indicated below.

\begin{table}[H]
    \centering
    \begin{tabular}{|p{4cm}|p{4cm}|l|}
        \hline
        \textbf{Suffix} & \textbf{Type} & \textbf{Examples} \\
        \hline
        None, only digits & \code{int} (\code{long} if too big) & \code{1234} \\
        \code{l} or \code{L} & \code{long} & \code
        {1234L} \\
        \code{u} or \code{U} & \code{unsigned int} (\code{unsigned long} if too big) & \code{1234U} \\
        \code{ul} or \code{UL} & \code{unsigned long} & \code{1234UL} \\
        None, with decimal or scientific & \code{double} & \code{1.234}, \code{1e-5} \\
        \code{f} or \code{F} & \code{float} & \code{1.234F} \\
        \code{l} or \code{L}, with decimal or scientific & \code{long double} & \code{1.234L}, \code{1e-5L} \\
        \hline
    \end{tabular}
    \caption{Data type and examples of integers using the given suffix}
\end{table}

\paragraph{Prefixes} Integer constants may contain the given prefixes below. These can also be used in conjunction with appropriate suffixes (those available for integers).

\begin{table}[H]
    \centering
    \begin{tabular}{|l|l|}
        \hline
        \textbf{Prefix} & \textbf{Interpretation} \\
        \hline
        None & Decimal \\
        \code{0} & Octal \\
        \code{0x} or \code{0X} & Hexadecimal \\
        \hline
    \end{tabular}
    \caption{Interpretations of integer constants with the given prefix}
\end{table}

\subsubsection{Character Constants}

\begin{itemize}
    \item Character constants are integers written as one character with single quotes and are equal to the numeric value in the machine's character set.
    \item Character constants can also participate in arithemtic operations like other integer types.
\end{itemize}

\paragraph{Escape Sequences} Certain characters can be represented by the following escape sequences.

\begin{table}[H]
    \centering
    \begin{tabular}{|l|l|l|}
        \hline
        \textbf{Sequence} & \textbf{Meaning} & \textbf{ASCII Code} \\
        \hline
        \code{\textbackslash 0} & Null character & 0 \\
        \code{\textbackslash a} & Alert (bell) character & 7 \\
        \code{\textbackslash b} & Backspace & 8 \\
        \code{\textbackslash t} & Horizontal tab & 9 \\
        \code{\textbackslash n} & Line feed (newline) & 10 \\
        \code{\textbackslash v} & Vertical tab & 11 \\
        \code{\textbackslash f} & Form feed & 12 \\
        \code{\textbackslash r} & Carriage return & 13 \\
        \code{\textbackslash "} & Double quote & 34 \\
        \code{\textbackslash '} & Single quote & 39 \\
        \code{\textbackslash ?} & Question mark & 63 \\
        \code{\textbackslash \textbackslash} & Backslash & 92 \\
        \code{\textbackslash ooo} & Octal code (o = octal digit) & \\
        \code{\textbackslash xhh} & Hex code (h = hex digit) & \\
        \hline
    \end{tabular}
    \caption{Different escape sequences, their meaning, and ASCII codes in C}
\end{table}

\subsubsection{Constant Expressions}

\begin{itemize}
    \item Constant expressions are expressions that involve only constants.
    \item Such expressions are typically evaluated at compile-time and can be used in any place that a constant can occur.
\end{itemize}

Example:

\begin{lstlisting}
#define LEAP 1;

int days[31+28+LEAP+31+30+31+30+31+31+30+31+30+31];
\end{lstlisting}

\subsubsection{String Constants}

\begin{itemize}
    \item A string constant (a.k.a. a string literal) is a sequence of 0 or more characters surrounded by double quotes.
    \item The escape sequences used in characters also apply in strings.
    \item String constants can be concatenated at compile-time, e.g \\\code{"hello, world"} is equivalent to \code{"hello ," " world"}.
    \item This is useful for breaking up long strings across several lines of code.
\end{itemize}

\subsubsection{Enum Constants}

\begin{itemize}
    \item An enum is a list of constant integer values.
    \item Names in enums \emph{must} be distinct.
    \item Values in enums don't need to be distinct.
    \item If values are left unassigned, they increase in order from the leftmost assigned value (the first will be 0 if unassigned).
\end{itemize}

Examples:

\begin{lstlisting}
enum boolean { NO, YES }; // NO=0, YES=1
enum letters { A = 10, B, C }; // A=10, B=11, C=12
enum letters2 { A, B = 50, C }; // A=0, B=50, C=51
\end{lstlisting}

\subsection{Declarations}

\begin{itemize}
    \item All variable declarations must be at the top of a code block before executable statements and declared before use.\footnote{Specific to C89/90.}
    \item A declaration contains a list of one or more variables of that type which may optionally be assigned.
    \item Automatic variables are local variables that are automatically allocated when the program enters the code block they are defined in.
    \item If a variable is not automatic, its initialization is done only once (before the program starts executing).
    \item Otherwise, it is initialised each time its code block is entered.
    \item External and \code{static} variables are initialised to 0 by default.
    \item Automatic variables with no initializer are undefined (garbage bits leftover in memory).
    \item The \code{const} qualifier can be applied to a variable declaration to specify its value will not be changed (for arrays the element will not be altered).
    \item \code{const} can also be used with function arguments to indicate that the function will not change, or reassign the local copy of the value passed.
\end{itemize}

Automatic variable examples:

\begin{lstlisting}
int a, b, c = 10; // a and b are undefined
const int d = 20; // will not be changed
\end{lstlisting}

\subsection{Operators}

\paragraph{Precedence} Operator precedence is defined according to the following table:

\begin{table}[H]
    \centering
    \begin{tabular}{|c|l|p{5cm}|c|}
        \hline
        \textbf{Rank} & \textbf{Category} & \textbf{Operators} & \textbf{Association} \\
        \hline
        1 & Postfix & \code{++}, \code{--}, \code{()}, \code{[]}, \code{->}, \code{.} & LTR \\
        \hline
        2 & Prefix & \code{+}, \code{-}, \code{++}, \code{--}, \code{!}, \code{~} & RTL \\
        \hline
        3 & Multiplicative & \code{*}, \code{/}, \code{\%} & LTR \\
        \hline
        4 & Additive & \code{+}, \code{-} & LTR \\
        \hline
        5 & Bit Shift & \code{<<}, \code{>>} & LTR \\
        \hline
        6 & Relational & \code{<}, \code{<=}, \code{>}, \code{>=} & LTR \\
        \hline
        7 & Equality & \code{==}, \code{!=} & LTR \\
        \hline
        8 & Bitwise AND & \code{\&} & LTR \\
        \hline
        9 & Bitwise XOR & \code{\^} & LTR \\
        \hline
        10 & Bitwise OR & \code{|} & LTR \\
        \hline
        11 & Logical AND & \code{\&\&} & LTR \\
        \hline
        12 & Logical OR & \code{||} & LTR \\
        \hline
        13 & Ternary & \code{?:} & RTL \\
        \hline
        14 & Assignment & \code{=}, \code{+=}, \code{-=}, \code{*=}, \code{/=}, \code{\%=}, \code{\&=}, \code{\^{}=}, \code{|=}, \code{<<=}, \code{>>=} & RTL \\
        \hline
        15 & Comma & \code{,} & LTR \\
        \hline
    \end{tabular}
    \caption{Operator precedence in C}
\end{table}

\paragraph{Arithmetic Operators}

\begin{itemize}
    \item Arithmetic operators consist of the additive and multiplicative operators.
    \item The \code{\%} operator cannot be applied to floats or doubles.
    \item The direction for truncation for \code{/} is towards zero.\footnote{In the C89/90 standard it is implementation dependent, can be towards zero or negative infinity but typically zero in modern standards.}
    \item The C standard requires that $(a / b)\cdot b + (a \% b) == a$ where $/$ is floor division.
    \item The sign of the result for \code{\%} usually takes on the sign of the dividend ($a$ for $a \% b$) as a result of \code{/} truncating towards zero.\footnote{For truncation towards negative infinity the sign of modulo results take on the sign of the divisor.}
\end{itemize}

\paragraph{Relational and Logical Operators}

\begin{itemize}
    \item The numerical value of expressions involving relational and logical operators is \code{1} if the expression is true or \code{0} if false.
\end{itemize}

\paragraph{Increment and Decrement Operators}

\begin{itemize}
    \item Prefix increment/decrement \emph{modifies the value then consumes}.
    \item Postfix increment/decrement \emph{consumes the value then modifies}.
    \item Increment/decrement operators can only be applied to variables (not expressions).
\end{itemize}

\paragraph{Bitwise Operators}

\begin{itemize}
    \item Bitwise operators may only be used on integer type operands.
    \item Bitwise AND is often used to mask off some bits (e.g \code{n \&= 0xF} sets all bits to \code{0} except the 4 lowest order bits).
    \item Bitwise OR is often to use to turn bits on (e.g \code{n |= 0xF} turns on the 4 lowest order bits).
    \item Shift operators \code{<<} and \code{>>} shift the bits up and down respectively of the left-hand operator by the non-negative amount in the right-hand operator (e.g \code{2 << 2} = \code{8} since \code{0010} becomes \code{1000}).
    \item Vacant bits after shifting are filled with 0s usually.
\end{itemize}

\paragraph{Ternary Operator}

\begin{itemize}
    \item The first expression of a ternary operator expression is evaluated.
    \item If the first is non-zero (true), the second is evaluated and that is the value of the ternary expression, otherwise (if it is false) the third expression is evaluated and \emph{that} is the value of the ternary expression.
    \item If the second and third expressions differ in type, then type conversion occurs based on the type conversion rules in the next section.
\end{itemize}

\subsection{Type Conversions}

\begin{itemize}
    \item Type conversions take place in things like expressions and assignments.
    \item Function prototypes enable automatic casting of arguments.
\end{itemize}

\subsubsection{Integer Promotion}

\begin{itemize}
    \item When a smaller rank is used in a binary arithmetic operation, it is automatically promoted to the larger rank used in the expression (promotion).
    \item So the operation can be performed as if it was on two operands of the same type.
    \item In an expression, type conversion happens incrementally along the way as the expression is evaluated according to associativity of each operator.
    \item \code{char}s and \code{short}s are always promoted to \code{int} (\code{signed}) in an expression.
\end{itemize}

\subsubsection{Order of Type Conversion Operations}

\begin{enumerate}
    \item \code{char}s and \code{short}s are promoted to \code{int} and any other integer promotion goes through.
    \item Signs are then compared using the size of the original types before promotion (\code{int} if originally \code{char} or \code{short}).
    \item If $a$ is \textbf{unsigned} and has size $\geq$ than $b$, a \textbf{signed} type, $b$ becomes unsigned.
    \item If $a$ is \textbf{unsigned} and has size < than $b$, a \textbf{signed} type, $a$ becomes signed.\footnote{The positive signed range of type $b$ usually matches the unsigned range of $a$, with the type of $b$ having twice the number of bits.}
\end{enumerate}

\subsubsection{Sign Extension and Zero Extension}

\begin{itemize}
    \item When a smaller type is unsigned and promoted, the program performs \emph{zero extension} i.e, fills the new high-order bits with 0s.
    \item When a smaller type is signed and promoted, the program performs \emph{sign extension} i.e, fills the new high-order bits with whatever the old two's complement sign bit was (0 if positive, 1 if negative).
    \item Each of these operations preserves the original value but the actual value depends on the new interpretation (i.e signed or unsigned).
\end{itemize}

\subsubsection{Truncation}

\begin{itemize}
    \item Larger integer types are converted to smaller ones by chopping off the excess higher order bits.
\end{itemize}

\subsubsection{Floats}

\begin{itemize}
    \item If integer types are mixed with float types, the float type always wins.
    \item If the integer type is larger than the float type, precision loss may occur.
\end{itemize}

\subsubsection{Casting}

\begin{itemize}
    \item An expression can be forcibly cast to another type.
    \item A cast behaves the same way as if an expression was being assigned to a variable of the cast type.
    \item Casting uses the following syntax:
    \begin{center}
        \textit{(type name) expression}
    \end{center}
\end{itemize}

\end{document}